\begin{lemma}[The unit-circle angle parametrization]
Let $x,y\in\mathbb{R}$ satisfy $x^2+y^2=1$. Then there exists
$\alpha\in\mathbb{R}$ with $\cos\alpha=x$ and $\sin\alpha=y$.
\end{lemma}
\begin{proof}
\emph{Step 1: $x$ lies in $[-1,1]$.} Since $y^2\ge 0$, the hypothesis gives
$x^2=1-y^2\le 1$. Expanding $(x+1)^2\ge 0$ gives $2x\ge -1-x^2$, and combining
with $x^2\le 1$ yields $2x\ge -2$, i.e.\ $-1\le x$. Expanding $(x-1)^2\ge 0$
gives $2x\le 1+x^2$, and combining with $x^2\le 1$ yields $2x\le 2$, i.e.\
$x\le 1$. Hence $\arccos x$ is defined and $\cos(\arccos x)=x$.

We also record that $1-x^2=y^2$, which is the hypothesis rearranged, and that
$\sin(\arccos x)=\sqrt{1-x^2}$ for $x\in[-1,1]$; so
$\sin(\arccos x)=\sqrt{y^2}$.

\emph{Step 2: the case $y\ge 0$.} Take $\alpha=\arccos x$. Then
$\cos\alpha=x$ by Step 1. For the sine, $\sin\alpha=\sqrt{y^2}$, and since
$y\ge 0$ we have $\sqrt{y^2}=y$. Hence $\sin\alpha=y$, and $\alpha$ is as
required.

\emph{Step 3: the case $y<0$.} Take $\alpha=-\arccos x$. The cosine is even, so
$\cos\alpha=\cos(\arccos x)=x$ by Step 1. The sine is odd, so
\[
  \sin\alpha=-\sin(\arccos x)=-\sqrt{1-x^2}=-\sqrt{y^2}=-\lvert y\rvert ,
\]
using $\sqrt{t^2}=\lvert t\rvert$ for every real $t$. Since $y<0$ we have
$\lvert y\rvert=-y$, so $\sin\alpha=-(-y)=y$. Hence $\alpha$ is again as
required.

Every real $y$ satisfies either $y\ge 0$ or $y<0$, so one of Steps 2 and 3
applies and the angle exists in all cases.
\end{proof}
